\documentclass[preview]{standalone}

\usepackage{amsmath}
\usepackage{amssymb}
\usepackage{tikz}
\usepackage{stellar}
\usepackage{definitions}
\usepackage{bettelini}

\begin{document}

\id{linearalgebra-bases}
\genpage

\section{Linear combinations}

\includesnpt{linear-combination-definition}

\begin{snippetexample}{linear-combination-r3-example}{Linear combinations in \(\mathbb{R}^3\)}
    In \(\realnumbers^3\), let
    \[
        v_1=(2,1,-3),
        \qquad
        v_2=(1,-2,4).
    \]
    Then \((17,-4,2)\) is a \linearcombination of \(v_1\) and \(v_2\),
    because
    \[
        (17,-4,2)=6(2,1,-3)+5(1,-2,4).
    \]
    On the other hand, \((17,-4,5)\) is not a \linearcombination of
    \(v_1\) and \(v_2\). If
    \[
        (17,-4,5)=x(2,1,-3)+y(1,-2,4),
    \]
    then
    \[
        \begin{cases}
            2x+y=17,\\
            x-2y=-4,\\
            -3x+4y=5.
        \end{cases}
    \]
    The first two equations give \(x=6\) and \(y=5\), but then
    \(-3x+4y=2\neq5\). Thus the system is inconsistent.
\end{snippetexample}

\section{Span}

\includesnpt{span-definition}

\begin{snippetproposition}{span-is-linear-subspace}{Span is the smallest generated subspace}
    Let \(V\) be a \vectorspace over a \field \(\mathbb{F}\), and let
    \(S\subseteq V\). Then \(\linearspan(S)\) is the smallest linear
    subspace of \(V\) containing every element of \(S\).
\end{snippetproposition}

\begin{snippetproof}{span-is-linear-subspace-proof}{span-is-linear-subspace}{Span is the smallest generated subspace}
    First, \(\linearspan(S)\) is a linear subspace. The zero vector belongs
    to it by the convention \(\linearspan(\emptyset)=\{0_V\}\) or by the
    trivial \linearcombination. If
    \[
        u=a_1v_1+\cdots+a_kv_k,
        \qquad
        w=b_1v_1+\cdots+b_kv_k
    \]
    are two \linearcombination[linear combinations] of vectors in \(S\), then
    \[
        u+w=(a_1+b_1)v_1+\cdots+(a_k+b_k)v_k
    \]
    is again a \linearcombination of vectors in \(S\). If
    \(\lambda\in\mathbb{F}\), then
    \[
        \lambda u=(\lambda a_1)v_1+\cdots+(\lambda a_k)v_k
    \]
    is again in \(\linearspan(S)\).

    Now let \(W\) be a linear subspace of \(V\) containing \(S\). Since \(W\)
    is closed under scalar multiplication and finite sums, it contains every
    finite \linearcombination of vectors in \(S\). Hence
    \(\linearspan(S)\subseteq W\). Therefore \(\linearspan(S)\) is the
    smallest linear subspace containing \(S\).
\end{snippetproof}

\begin{snippetdefinition}{vector-space-generating-set}{Generating set}
    Let \(V\) be a \vectorspace over a \field \(\mathbb{F}\). A finite
    collection of \vector[vectors]
    \[
        \{v_1,\ldots,v_m\}\subseteq V
    \]
    is a \emph{generating set} for \(V\) if
    \[
        \linearspan(v_1,\ldots,v_m)=V.
    \]
    Equivalently, every \vector of \(V\) is a \linearcombination of
    \(v_1,\ldots,v_m\).
\end{snippetdefinition}

\begin{snippetexample}{matrix-space-standard-generators-example}{Generators of \(\matrices_{2\times2}(\mathbb{R})\)}
    Let \(V=\matrices_{2\times2}(\realnumbers)\), and let
    \[
        E_{11}=
        \begin{pmatrix}
            1 & 0\\
            0 & 0
        \end{pmatrix},
        \quad
        E_{12}=
        \begin{pmatrix}
            0 & 1\\
            0 & 0
        \end{pmatrix},
    \]
    \[
        E_{21}=
        \begin{pmatrix}
            0 & 0\\
            1 & 0
        \end{pmatrix},
        \quad
        E_{22}=
        \begin{pmatrix}
            0 & 0\\
            0 & 1
        \end{pmatrix}.
    \]
    Then \(\{E_{11},E_{12},E_{21},E_{22}\}\) generates \(V\), because every
    \matrix
    \[
        A=
        \begin{pmatrix}
            x & y\\
            z & t
        \end{pmatrix}
    \]
    can be written as
    \[
        A=xE_{11}+yE_{12}+zE_{21}+tE_{22}.
    \]
\end{snippetexample}

\section{Finite generation}

\begin{snippetdefinition}{finite-vector-space-definition}{Finite-dimensional vector space}
    A \vectorspace \(V\) over a \field \(\mathbb{F}\) is
    \emph{finite-dimensional} if \(V\) admits a finite generating \set.
\end{snippetdefinition}

\begin{snippetexample}{fn-finite-dimensional-example}{The space \(F^n\) is finite-dimensional}
    Let \(\mathbb{F}\) be a \field. The \vectorspace \(\mathbb{F}^n\) is
    finite-dimensional. If
    \[
        e_1=(1,0,\ldots,0),
        \quad
        e_2=(0,1,0,\ldots,0),
        \quad
        \ldots,
        \quad
        e_n=(0,\ldots,0,1),
    \]
    then every \(x=(x_1,\ldots,x_n)\in\mathbb{F}^n\) satisfies
    \[
        x=x_1e_1+\cdots+x_ne_n.
    \]
    Thus \(\{e_1,\ldots,e_n\}\) generates \(\mathbb{F}^n\).
\end{snippetexample}

\begin{snippetexample}{matrix-space-finite-dimensional-example}{Matrix spaces are finite-dimensional}
    The \vectorspace \(\matrices_{m\times n}(\mathbb{F})\) is
    finite-dimensional. For \(1\leq i\leq m\) and \(1\leq j\leq n\), let
    \(E_{ij}\) be the \matrix with \(1\) in position \((i,j)\) and \(0\) in
    every other position. Then the finite collection
    \[
        \{E_{ij}\suchthat 1\leq i\leq m,\ 1\leq j\leq n\}
    \]
    generates \(\matrices_{m\times n}(\mathbb{F})\), since
    \[
        A=(a_{ij})
        =
        \sum_{i=1}^m\sum_{j=1}^n a_{ij}E_{ij}.
    \]
\end{snippetexample}

\begin{snippetdefinition}{infinite-dimensional-vector-space-definition}{Infinite-dimensional vector space}
    A \vectorspace is \emph{infinite-dimensional} if it is not
    finite-dimensional, that is, if it admits no finite generating \set.
\end{snippetdefinition}

\begin{snippetexample}{bounded-degree-polynomials-finite-dimensional-example}{Polynomials of bounded degree}
    Let
    \[
        \mathcal{P}_m(\mathbb{F})
        =
        \{p\in\mathbb{F}[x]\suchthat \deg(p)\leq m\}\union\{0\}.
    \]
    Then \(\mathcal{P}_m(\mathbb{F})\) is a linear subspace of
    \(\mathbb{F}[x]\), and it is finite-dimensional because
    \[
        \{1,x,x^2,\ldots,x^m\}
    \]
    generates it.
\end{snippetexample}

\begin{snippetexample}{polynomial-space-infinite-dimensional-example}{The polynomial space is infinite-dimensional}
    The \vectorspace \(\mathbb{F}[x]\) is infinite-dimensional. Let
    \(\{p_1,\ldots,p_m\}\) be any finite collection of \polynomial[polynomials]. If
    \(d_i=\deg(p_i)\) and \(d=\max\{d_1,\ldots,d_m\}\), then every linear
    combination of \(p_1,\ldots,p_m\) has degree at most \(d\). Hence
    \(x^{d+1}\) is not in their span. No finite collection can generate
    \(\mathbb{F}[x]\).
\end{snippetexample}

\section{Linear independence}

\includesnpt{linear-independence-definition}

\begin{snippetdefinition}{linear-dependence-definition}{Linear dependence}
    Let \(V\) be a \vectorspace over a \field \(\mathbb{F}\). A finite
    collection
    \[
        \{v_1,\ldots,v_m\}\subseteq V
    \]
    is \emph{linearly dependent} if it is not \linearlyindependent.
    Equivalently, there exist \vsscalar[scalars] \(a_1,\ldots,a_m\in\mathbb{F}\), not
    all zero, such that
    \[
        a_1v_1+\cdots+a_mv_m=0_V.
    \]
\end{snippetdefinition}

\begin{snippetproposition}{linear-independent-coordinates-unique}{Uniqueness of coordinates}
    Let \(V\) be a \vectorspace over a \field \(\mathbb{F}\), and let
    \(v_1,\ldots,v_m\in V\). For every
    \(v\in\linearspan(v_1,\ldots,v_m)\), the expression
    \[
        v=a_1v_1+\cdots+a_mv_m
    \]
    is unique if and only if \(\{v_1,\ldots,v_m\}\) is \linearlyindependent.
\end{snippetproposition}

\begin{snippetproof}{linear-independent-coordinates-unique-proof}{linear-independent-coordinates-unique}{Uniqueness of coordinates}
    Suppose first that \(\{v_1,\ldots,v_m\}\) is \linearlyindependent. If
    \[
        v=a_1v_1+\cdots+a_mv_m=b_1v_1+\cdots+b_mv_m,
    \]
    then
    \[
        (a_1-b_1)v_1+\cdots+(a_m-b_m)v_m=0_V.
    \]
    Linear independence gives \(a_i-b_i=0\) for every \(i\), so \(a_i=b_i\).

    Conversely, suppose that every \vector in the span has a unique
    expression. The zero \vector has the expression
    \[
        0_V=0v_1+\cdots+0v_m.
    \]
    If there were a nontrivial linear relation
    \[
        a_1v_1+\cdots+a_mv_m=0_V,
    \]
    then \(0_V\) would have two different expressions. Hence every such
    relation is trivial, and the collection is \linearlyindependent.
\end{snippetproof}

\begin{snippetexample}{single-vector-linear-independence-example}{One vector}
    A collection \(\{v\}\) is \linearlyindependent if and only if
    \(v\neq0_V\). If \(v\neq0_V\) and \(av=0_V\), then \(a=0\). If
    \(v=0_V\), then \(1v=0_V\) is a nontrivial linear relation.
\end{snippetexample}

\begin{snippetexample}{two-vector-linear-independence-example}{Two vectors}
    A collection \(\{v,w\}\) is \linearlyindependent if and only if neither
    \vector is a scalar multiple of the other. If \(v=\lambda w\), then
    \(v-\lambda w=0_V\) is a nontrivial relation. Conversely, if
    \(av+bw=0_V\) with \(a\neq0\), then \(v=-\frac{b}{a}w\); if \(b\neq0\),
    then \(w=-\frac{a}{b}v\).
\end{snippetexample}

\begin{snippetexample}{canonical-basis-linearly-independent-example}{Independence of the canonical vectors}
    The collection \(\{e_1,\ldots,e_n\}\) is \linearlyindependent in
    \(\mathbb{F}^n\). If
    \[
        a_1e_1+\cdots+a_ne_n=0,
    \]
    then
    \[
        (a_1,\ldots,a_n)=(0,\ldots,0),
    \]
    so \(a_1=\cdots=a_n=0\).
\end{snippetexample}

\begin{snippetexample}{polynomial-monomials-linearly-independent-example}{Polynomial monomials}
    The collection
    \[
        \{1,x,x^2,\ldots,x^m\}
    \]
    is \linearlyindependent in \(\mathcal{P}_m(\mathbb{F})\). If
    \[
        a_0+a_1x+\cdots+a_mx^m=0
    \]
    as a \polynomial, then all coefficients are zero.
\end{snippetexample}

\begin{snippetproposition}{linearly-independent-subcollection}{Subcollections of independent collections}
    Every subcollection of a \linearlyindependent collection is
    \linearlyindependent.
\end{snippetproposition}

\begin{snippetproof}{linearly-independent-subcollection-proof}{linearly-independent-subcollection}{Subcollections of independent collections}
    Let \(\{v_1,\ldots,v_m\}\) be \linearlyindependent, and consider a
    subcollection \(\{v_{i_1},\ldots,v_{i_k}\}\). Any linear relation among
    the \(v_{i_j}\) can be viewed as a linear relation among all
    \(v_1,\ldots,v_m\) by assigning coefficient \(0\) to the omitted \vector[vectors].
    Since the larger collection is \linearlyindependent, all coefficients in
    the relation are zero.
\end{snippetproof}

\begin{snippetexample}{dependent-vectors-r3-example}{Three dependent vectors in \(\mathbb{R}^3\)}
    In \(\realnumbers^3\), let
    \[
        v_1=(2,3,1),
        \qquad
        v_2=(1,-1,2),
        \qquad
        v_3=(7,3,8).
    \]
    Then the collection is linearly dependent because
    \[
        2v_1+3v_2-v_3=(0,0,0).
    \]
\end{snippetexample}

\begin{snippetexample}{f2-linear-dependence-example}{Dependence over \(F_2\)}
    In \(F_2^3\), let
    \[
        v_1=(1,1,0),
        \qquad
        v_2=(1,0,1),
        \qquad
        v_3=(0,1,1).
    \]
    Then
    \[
        v_1+v_2+v_3=(0,0,0),
    \]
    because \(1+1=0\) in \(F_2\). Thus the collection is linearly dependent
    over \(F_2\). The same collection, viewed in \(\realnumbers^3\), is
    \linearlyindependent.
\end{snippetexample}

\begin{snippetproposition}{vector-linear-combination-of-others-dependent}{A redundant vector implies dependence}
    Let \(V\) be a \vectorspace over a \field \(\mathbb{F}\). If one \vector in
    a finite collection is a \linearcombination of the others, then the
    collection is linearly dependent.
\end{snippetproposition}

\begin{snippetproof}{vector-linear-combination-of-others-dependent-proof}{vector-linear-combination-of-others-dependent}{A redundant vector implies dependence}
    Suppose, after reordering if necessary, that
    \[
        v_1=a_2v_2+\cdots+a_mv_m.
    \]
    Then
    \[
        v_1-a_2v_2-\cdots-a_mv_m=0_V
    \]
    is a nontrivial linear relation, since the coefficient of \(v_1\) is
    \(1\).
\end{snippetproof}

\begin{snippetlemma}{linearly-dependent-set-reduction}{Linear dependence lemma}
    Let \(V\) be a \vectorspace over a \field \(\mathbb{F}\), and let
    \[
        \{v_1,\ldots,v_m\}
    \]
    be an ordered linearly dependent collection of \vector[vectors]. Then there exists
    an index \(j\) such that
    \[
        v_j\in\linearspan(v_1,\ldots,v_{j-1})
    \]
    and
    \[
        \linearspan(v_1,\ldots,v_{j-1},v_{j+1},\ldots,v_m)
        =
        \linearspan(v_1,\ldots,v_m).
    \]
\end{snippetlemma}

\begin{snippetproof}{linearly-dependent-set-reduction-proof}{linearly-dependent-set-reduction}{Linear dependence lemma}
    Since the collection is linearly dependent, there are \vsscalar[scalars]
    \(a_1,\ldots,a_m\in\mathbb{F}\), not all zero, such that
    \[
        a_1v_1+\cdots+a_mv_m=0_V.
    \]
    Let \(j\) be the largest index with \(a_j\neq0\). Then
    \[
        a_1v_1+\cdots+a_jv_j=0_V,
    \]
    so
    \[
        v_j
        =
        -\frac{a_1}{a_j}v_1-\cdots-\frac{a_{j-1}}{a_j}v_{j-1}.
    \]
    Hence \(v_j\in\linearspan(v_1,\ldots,v_{j-1})\). Removing \(v_j\)
    cannot enlarge the span. Conversely, since \(v_j\) is already in the
    span of the remaining vectors, every \linearcombination using \(v_j\) can
    be rewritten without \(v_j\). Therefore the spans are equal.
\end{snippetproof}

\section{Bases}

\begin{snippetdefinition}{basis-definition}{Basis}
    Let \(V\) be a \vectorspace over a \field \(\mathbb{F}\). A
    \emph{basis} of \(V\) is a finite collection
    \[
        B=\{v_1,\ldots,v_n\}\subseteq V
    \]
    that generates \(V\) and is \linearlyindependent.
\end{snippetdefinition}

\begin{snippet}{basis-size-and-redundancy}
    A \basis is large enough to generate the whole \vectorspace and small
    enough not to contain redundant \vector[vectors].
\end{snippet}

\begin{snippetexample}{canonical-basis-fn-example}{Canonical basis of \(F^n\)}
    For \(j=1,\ldots,n\), let
    \[
        e_j=(0,\ldots,0,1,0,\ldots,0),
    \]
    where \(1\) is in the \(j\)-th position. Then
    \[
        \{e_1,\ldots,e_n\}
    \]
    is a \basis of \(\mathbb{F}^n\), because every
    \(x=(x_1,\ldots,x_n)\) has the unique expression
    \[
        x=x_1e_1+\cdots+x_ne_n.
    \]
\end{snippetexample}

\begin{snippetexample}{canonical-basis-matrix-space-example}{Canonical basis of a matrix space}
    The \matrix[matrices]
    \[
        \{E_{ij}\suchthat 1\leq i\leq m,\ 1\leq j\leq n\},
    \]
    where \(E_{ij}\) has \(1\) in position \((i,j)\) and \(0\) elsewhere,
    form a \basis of \(\matrices_{m\times n}(\mathbb{F})\). Indeed,
    \[
        A=(a_{ij})=\sum_{i=1}^m\sum_{j=1}^n a_{ij}E_{ij},
    \]
    and this expression is unique.
\end{snippetexample}

\begin{snippetexample}{polynomial-bounded-degree-basis-example}{Basis of bounded-degree polynomials}
    The \vectorspace
    \[
        \mathcal{P}_m(\mathbb{F})
        =
        \{p\in\mathbb{F}[x]\suchthat\deg(p)\leq m\}\union\{0\}
    \]
    has \basis
    \[
        \{1,x,x^2,\ldots,x^m\}.
    \]
\end{snippetexample}

\begin{snippetexample}{basis-of-r2-example}{Bases of \(\mathbb{R}^2\)}
    The collection \(\{(1,2),(3,4)\}\) is a \basis of \(\realnumbers^2\).
    For a \vector \((b_1,b_2)\), solving
    \[
        x_1(1,2)+x_2(3,4)=(b_1,b_2)
    \]
    is equivalent to
    \[
        \begin{cases}
            x_1+3x_2=b_1,\\
            2x_1+4x_2=b_2.
        \end{cases}
    \]
    The coefficient determinant is \(1\cdot4-3\cdot2=-2\neq0\), so the
    system has a unique solution for every \((b_1,b_2)\).
\end{snippetexample}

\begin{snippetexample}{non-bases-examples}{Collections that are not bases}
    The collection
    \[
        \{(1,2),(3,5),(4,13)\}
    \]
    is not a \basis of \(\realnumbers^2\), because it is linearly dependent:
    \[
        19(1,2)-5(3,5)-(4,13)=0.
    \]
    The collection
    \[
        \{(1,2,-4),(7,-5,6)\}
    \]
    is not a \basis of \(\realnumbers^3\), because two \vector[vectors] cannot generate
    all of \(\realnumbers^3\).
\end{snippetexample}

\begin{snippetexample}{trace-zero-matrices-basis-example}{Basis of trace-zero matrices}
    Let
    \[
        W=
        \left\{
            \begin{pmatrix}
                a & b\\
                c & d
            \end{pmatrix}
            \in\matrices_{2\times2}(\realnumbers)
            \suchthat a+d=0
        \right\}.
    \]
    Then
    \[
        \left\{
            \begin{pmatrix}
                1 & 0\\
                0 & -1
            \end{pmatrix},
            \begin{pmatrix}
                0 & 1\\
                0 & 0
            \end{pmatrix},
            \begin{pmatrix}
                0 & 0\\
                1 & 0
            \end{pmatrix}
        \right\}
    \]
    is a \basis of \(W\), since every element of \(W\) has the unique form
    \[
        \begin{pmatrix}
            a & b\\
            c & -a
        \end{pmatrix}
        =
        a
        \begin{pmatrix}
            1 & 0\\
            0 & -1
        \end{pmatrix}
        +
        b
        \begin{pmatrix}
            0 & 1\\
            0 & 0
        \end{pmatrix}
        +
        c
        \begin{pmatrix}
            0 & 0\\
            1 & 0
        \end{pmatrix}.
    \]
\end{snippetexample}

\begin{snippetproposition}{vector-space-basis-criterion}{Basis criterion}
    Let \(V\) be a \vectorspace over a \field \(\mathbb{F}\), and let
    \(B=\{v_1,\ldots,v_n\}\subseteq V\). Then \(B\) is a \basis of \(V\) if
    and only if every \(v\in V\) can be written uniquely as
    \[
        v=a_1v_1+\cdots+a_nv_n.
    \]
\end{snippetproposition}

\begin{snippetproof}{vector-space-basis-criterion-proof}{vector-space-basis-criterion}{Basis criterion}
    \begin{iffproof}
        \item If \(B\) is a \basis, then it generates \(V\), so every \vector
        has at least one expression as a \linearcombination of \(B\). Since
        \(B\) is \linearlyindependent, the uniqueness proposition for
        coordinates gives uniqueness.

        \item If every \vector of \(V\) has a unique expression as a
        \linearcombination of \(B\), then \(B\) generates \(V\). The zero \vector has
        the expression
        \[
            0_V=0v_1+\cdots+0v_n.
        \]
        By uniqueness, any relation
        \(a_1v_1+\cdots+a_nv_n=0_V\) has \(a_1=\cdots=a_n=0\). Hence \(B\)
        is \linearlyindependent, and therefore it is a \basis.
    \end{iffproof}
\end{snippetproof}

\section{Exchange principles}

\begin{snippettheorem}{vector-space-generative-set-cardinality-inequality}{Independent collections are no larger than generating collections}
    Let \(V\) be a finite-dimensional \vectorspace over a \field
    \(\mathbb{F}\). If \(\{v_1,\ldots,v_m\}\) is \linearlyindependent and
    \(\{w_1,\ldots,w_n\}\) generates \(V\), then
    \[
        m\leq n.
    \]
\end{snippettheorem}

\begin{snippetproof}{vector-space-generative-set-cardinality-inequality-proof}{vector-space-generative-set-cardinality-inequality}{Independent collections are no larger than generating collections}
    Starting from the generating collection \(\{w_1,\ldots,w_n\}\), add
    \(v_1\). Since the \(w_i\) generate \(V\), the enlarged collection is
    linearly dependent. By the linear dependence lemma, one of the \(w_i\)
    can be removed without changing the span; \(v_1\) cannot be removed,
    because it is nonzero.

    Repeat the same operation with \(v_2,\ldots,v_m\). After \(r\) steps,
    one obtains a generating collection of \(n\) \vector[vectors] containing
    \(v_1,\ldots,v_r\). If \(m>n\), after \(n\) steps the \vector \(v_{n+1}\)
    would be in the span of \(v_1,\ldots,v_n\), contradicting linear
    independence. Hence \(m\leq n\).
\end{snippetproof}

\begin{snippetcorollary}{finite-vector-space-subspace-is-finite}{Subspaces of finite-dimensional spaces}
    Every linear subspace of a finite-dimensional \vectorspace is
    finite-dimensional.
\end{snippetcorollary}

\begin{snippetproof}{finite-vector-space-subspace-is-finite-proof}{finite-vector-space-subspace-is-finite}{Subspaces of finite-dimensional spaces}
    Let \(V\) be finite-dimensional and let \(U\subseteq V\) be a linear
    subspace. If \(U=\{0_V\}\), then \(U\) is generated by the empty
    collection. Otherwise, choose \(u_1\in U\setminus\{0_V\}\). If \(U\) is
    not generated by \(u_1\), choose \(u_2\in U\setminus\linearspan(u_1)\),
    and continue. At every step the chosen \vector[vectors] are \linearlyindependent.
    This process cannot continue beyond the size of a finite generating
    collection of \(V\), by the previous theorem. Therefore it stops after
    finitely many steps and yields a finite generating collection for \(U\).
\end{snippetproof}

\begin{snippetexample}{sequence-space-infinite-dimensional-example}{The sequence space is infinite-dimensional}
    The \vectorspace \(\mathbb{F}^{\naturalnumbers}\) is
    infinite-dimensional. For every \(m\in\naturalnumbers\), let
    \(e_1,\ldots,e_m\) be the sequences where \(e_i\) is \(1\) in the
    \(i\)-th position and \(0\) elsewhere. Then
    \[
        \{e_1,\ldots,e_m\}
    \]
    is \linearlyindependent. Since such independent collections can have
    arbitrarily large finite size, no finite collection can generate
    \(\mathbb{F}^{\naturalnumbers}\).
\end{snippetexample}

\begin{snippetproposition}{vector-space-generating-set-reduction}{Reduction of a generating collection}
    Every finite generating collection of a \vectorspace can be reduced to a
    \basis.
\end{snippetproposition}

\begin{snippetproof}{vector-space-generating-set-reduction-proof}{vector-space-generating-set-reduction}{Reduction of a generating collection}
    Let \(\{v_1,\ldots,v_n\}\) generate \(V\). Inspect the \vector[vectors] in order.
    Remove \(v_1\) if \(v_1=0_V\). At step \(j\), remove \(v_j\) if it lies
    in the span of the previously retained \vector[vectors]; otherwise retain it.
    Removing such a \vector does not change the span. After finitely many
    steps, the retained collection still generates \(V\). It is linearly
    independent, because no retained \vector lies in the span of the previous
    retained \vector[vectors]. Hence it is a \basis.
\end{snippetproof}

\begin{snippetexample}{generating-set-reduction-r2-example}{Reducing generators in \(\mathbb{R}^2\)}
    The collection
    \[
        \{(1,2),(3,5),(4,7),(5,9)\}
    \]
    generates \(\realnumbers^2\). The \vector[vectors] \((4,7)\) and \((5,9)\) lie
    in the span of \((1,2)\) and \((3,5)\), while \((3,5)\) is not in the
    span of \((1,2)\). Reducing the generating collection gives the \basis
    \[
        \{(1,2),(3,5)\}.
    \]
\end{snippetexample}

\begin{snippetcorollary}{every-finite-vector-space-has-basis}{Existence of bases}
    Every finite-dimensional \vectorspace admits a \basis.
\end{snippetcorollary}

\begin{snippetproof}{every-finite-vector-space-has-basis-proof}{every-finite-vector-space-has-basis}{Existence of bases}
    By definition, a finite-dimensional \vectorspace admits a finite
    generating collection. By the reduction theorem, this collection can be
    reduced to a \basis.
\end{snippetproof}

\begin{snippetproposition}{linearly-independent-set-basis-expansion}{Basis extension theorem}
    Let \(V\) be a finite-dimensional \vectorspace over a \field
    \(\mathbb{F}\). Every \linearlyindependent finite collection in \(V\) can
    be extended to a \basis of \(V\).
\end{snippetproposition}

\begin{snippetproof}{linearly-independent-set-basis-expansion-proof}{linearly-independent-set-basis-expansion}{Basis extension theorem}
    Let \(\{u_1,\ldots,u_m\}\) be \linearlyindependent. Since \(V\) is
    finite-dimensional, it has a \basis \(\{w_1,\ldots,w_n\}\). The collection
    \[
        \{u_1,\ldots,u_m,w_1,\ldots,w_n\}
    \]
    generates \(V\). Reduce it to a \basis by removing \vector[vectors] that lie in the
    span of previous retained \vector[vectors]. None of the \(u_i\) is removed,
    because they are \linearlyindependent. The final \basis therefore contains
    \(u_1,\ldots,u_m\).
\end{snippetproof}

\begin{snippetexample}{basis-extension-r3-example}{Extending a collection in \(\mathbb{R}^3\)}
    The collection
    \[
        \{(2,3,4),(9,6,8)\}
    \]
    is \linearlyindependent in \(\realnumbers^3\). Starting from the
    canonical \basis, one checks that \(e_1=(1,0,0)\) belongs to the span of
    the first two \vector[vectors], while \(e_2=(0,1,0)\) does not. Therefore
    \[
        \{(2,3,4),(9,6,8),(0,1,0)\}
    \]
    is a \basis of \(\realnumbers^3\).
\end{snippetexample}

\section{Complements}

\begin{snippetdefinition}{linear-subspace-complement-definition}{Linear subspace complement}
    Let \(V\) be a finite-dimensional \vectorspace over a \field
    \(\mathbb{F}\), and let \(W\subseteq V\) be a linear subspace. A
    \emph{complement} of \(W\) in \(V\) is a linear subspace
    \(U\subseteq V\) such that
    \[
        V=W\oplus U.
    \]
    Equivalently, \(W+U=V\) and \(W\intersection U=\{0_V\}\).
\end{snippetdefinition}

\begin{snippetproposition}{linear-subspace-complement-existance}{Existence of complements}
    Let \(V\) be a finite-dimensional \vectorspace over a \field
    \(\mathbb{F}\), and let \(W\subseteq V\) be a linear subspace. Then
    \(W\) admits a complement in \(V\).
\end{snippetproposition}

\begin{snippetproof}{linear-subspace-complement-existance-proof}{linear-subspace-complement-existance}{Existence of complements}
    Since \(W\) is a subspace of a finite-dimensional \vectorspace, \(W\) is
    finite-dimensional. Let \(\{w_1,\ldots,w_m\}\) be a \basis of \(W\).
    Extend it to a \basis of \(V\):
    \[
        \{w_1,\ldots,w_m,u_1,\ldots,u_n\}.
    \]
    Put \(U=\linearspan(u_1,\ldots,u_n)\). Every \(v\in V\) can be written
    as
    \[
        v=a_1w_1+\cdots+a_mw_m+b_1u_1+\cdots+b_nu_n,
    \]
    where the first part belongs to \(W\) and the second belongs to \(U\).
    Hence \(W+U=V\).

    If \(v\in W\intersection U\), then
    \[
        v=a_1w_1+\cdots+a_mw_m=b_1u_1+\cdots+b_nu_n.
    \]
    Subtracting gives a linear relation among the \basis \vector[vectors]
    \(w_1,\ldots,w_m,u_1,\ldots,u_n\). All coefficients are zero, so
    \(v=0_V\). Thus \(W\intersection U=\{0_V\}\), and \(V=W\oplus U\).
\end{snippetproof}

\begin{snippetexample}{nonunique-complements-r2-example}{Complements are not unique}
    Let \(V=\realnumbers^2\) and \(W=\linearspan((1,0))\). For every
    \(k\in\realnumbers\), put
    \[
        U_k=\linearspan((k,1)).
    \]
    Then \(U_k\) is a complement of \(W\). Indeed, every
    \((x,y)\in\realnumbers^2\) can be written as
    \[
        (x,y)=(x-ky)(1,0)+y(k,1),
    \]
    and \(W\intersection U_k=\{0\}\). Hence
    \[
        \realnumbers^2=W\oplus U_k.
    \]
    \[
        \begin{tikzpicture}[scale=0.9]
            \draw[->] (-2,0) -- (3,0) node[right] {\(W\)};
            \draw[->] (0,-1.5) -- (0,2.5);
            \draw[->, thick] (0,0) -- (1.4,0) node[below] {\((1,0)\)};
            \draw[->] (0,0) -- (0.8,1.6) node[right] {\(U_k\)};
            \draw[->] (0,0) -- (-1.1,1.5) node[left] {\(U_h\)};
            \draw[->] (0,0) -- (0,1.8) node[left] {\(U_0\)};
            \fill (0,0) circle (1.3pt);
        \end{tikzpicture}
    \]
\end{snippetexample}

\end{document}
